% DEFINED:   sec:interp-setup, sec:circuit-restated
% RESOLVED:  ch:transformer, sec:one-layer-limit
% FORWARD:   (none)

\section{The setup}\label{sec:interp-setup}

The model under the lens is the one \Cref{ch:transformer} trained: a two-layer,
one-head decoder transformer with $d_{\text{model}} = 32$ and
$d_{\text{ff}} = 64$, trained for 2000 iterations with Adam at learning rate
$10^{-3}$, batch size 32, on 20{,}000 examples of the $M = 8$, $A = 3$ pointer
task. On 2000 held-out examples it reaches a test accuracy of 1.000. It has
17{,}282 parameters, and the task has only $2^{11} = 2048$ possible inputs. That
combination is what makes the dissection clean. The model solves the task
perfectly, and the task is small enough that we can describe what every
component does and do it from the trained weights and attention patterns
directly.

Each input is an eleven-token sequence: eight memory bits, then three address
bits, most significant first. The supervised target at position 10, the last
input position, is the value at memory position $a$, where $a$ is the integer
decoded from the three address bits. The plan for the chapter is to read the
architectural prediction of \Cref{ch:transformer} off these trained weights and
check, component by component, whether the model does what the architecture said
it must.

The instrument is the attention pattern itself. Both layers cache their
forward-pass attention matrix, and the model exposes a method that returns the
$(T, T)$ attention pattern for any single layer and head. We call it to extract
the pattern for a given input, and we average it across the test set to ask what
the model does in general rather than on one example. All of the code for this
analysis is in \texttt{examples/pointer\_interp.py}.

\section{What the circuit must be}\label{sec:circuit-restated}

Before looking at a single weight, it is worth restating what the model
\emph{has} to be doing, because the architectural argument of
\Cref{sec:one-layer-limit} is specific enough to pin the circuit down to one
possibility.

The model produces its prediction by passing the residual stream at position 10
through the final layer norm and an output head. Whatever the model knows about
the answer has to be present in that residual stream at position 10 after the
second layer. Now trace backward. Before any attention runs, position 10's
residual is the embedding of token $a_2$, the address least significant bit,
plus the positional encoding for position 10. It carries no information about
$a_0$ or $a_1$, which live in positions 8 and 9. Yet the answer depends on the
full address $a = 4 a_0 + 2 a_1 + a_2$. So for the second layer's query at
position 10 to attend to memory position $a$, that query must already depend on
all three address bits, which means the first layer must have moved $a_0$ and
$a_1$ into position 10's residual. The only mechanism a self-attention layer has
for that is attention from position 10 back to positions 8 and 9.

The hypothesis is therefore forced:

\begin{quote}
\textbf{Layer 1}, at position 10, attends to positions 8, 9, and 10, the three
address positions, and assembles them into position 10's residual.

\textbf{Layer 2}, at position 10, uses the assembled address to attend to memory
position $a$ and copy its value into the readout.
\end{quote}

This is the \emph{only} way a two-layer, one-head transformer can solve the task.
The claim is genuinely falsifiable. If the trained model is doing something else,
the architectural argument of \Cref{ch:transformer} was wrong. If it is doing
exactly this, the argument is empirically vindicated. The next three sections
check each half against the weights and then confirm the whole causally.

A note on scope before we begin. The model dissected here, and through the
induction-head and in-context-learning discussion that follows, is this
from-scratch NumPy model at $M = 8$. The final section turns the same tools on a
larger PyTorch model at $M = 32$, where the answer is more surprising.
